\section{Problem Formulation and Theoretical Analysis}
\label{sec:analysis}
\begin{comment}
visuomotor imitation learning
任务定义是这样定义，但是要强调我们的成立前提是这个任务的结果是可以被视觉评估

有限机器人示范下的离线视觉运动模仿学习，面向完成条件可通过视觉历史评估的操作任务。
环境信息由视觉提供，允许使用机器人本体感知与因果执行历史，不依赖物体特权状态或力触觉测量。
训练使用预先采集的固定示范集；部署期间不查询专家、不更新网络参数，但通过新观测和历史进行闭环决策。
在声明的相机配置、时间范围与可见性条件下，视觉历史能够提供判定任务完成的充分证据。

- 推动、摆放、分拣和可见区域内的物体整理。
- 抓取后搬运、放置与具有可见结果的堆叠。
- 部分布料展开、覆盖、折叠和工具操作。

> 在固定且有限的机器人示范条件下，对于环境信息由视觉提供、任务完成可由视觉历史评估的操作任务，如何通过语义引导的任务关系表示与历史条件反馈，提高离线训练的动作块模仿策略的任务完成表
> 现？


> 策略允许使用本体感知与因果执行历史，部署时进行闭环推理，但不依赖在线专家查询或参数更新。

在一切开始之前，先想清楚，要研究的问题是什么。
我觉得需要限定进来，因为我们研究的并不是一般的模仿学习

1 DAgger的想法是出现错误机器人进一步执行，会改变整个过程，所以需要处理机器人没遇到的状态迭代，进行在线学习对吧
2 Tu 等人的论文等人的想法是上面的进一步分析，有些影响会逐渐变大，有些影响是会衰减的，增量稳定性的概念，初始状态差异，输入扰动，对后续的轨迹的影响
（是否有给出具体的条件，怎么判别和计算？）
模仿误差进入专家闭环，能不能恢复回来
3 TaSIL不仅学习专家动作，还学习专家的局部变化规律，动作对状态的导数，可以约束附近的行为

既有研究揭示，模仿学习的闭环表现既受到策略诱导分布偏移的影响，也受到闭环系统误差传播特性的影响。
本文不假定语义表示能够补全未覆盖状态，也不假定历史反馈自动满足稳定性约束，
而是研究任务关系监督与计划条件反馈能否在有限示范设置下改善关键执行误差及任务完成表现。

% 本章服务于同一研究问题：有限示范下，任务关系表示与历史条件反馈如何支持任务完成表现的改善。
% 分析不将每个模块作为独立研究目标，而是解释两条互补路径。

% 首先简述动作误差与任务完成之间的条件性联系，说明不能仅以训练损失评判方法；
% 随后分析语义与几何监督如何缓解任务相关表示不足，
% 以及新观测和历史如何提供针对当前计划的增量纠偏信息。

% 覆盖界提供适用边界，表示分解解释监督设计，残差对齐与时序条件指导反馈验证。
% 以下命题不假定加入语义、查询或记忆就必然提高成功率，
% 也不以已知恒等式替代完整系统的实验证据。

\end{comment}

% 核心观点：研究有限固定示范下的离线视觉动作模仿学习，以视觉可评估的任务完成为部署目标。
% 叙述逻辑：系统与观测 -> 固定示范和因果策略 -> 区分训练代理与任务目标 -> 视觉可评估性假设。
% 边界：终局评估可辨识性不推出中间决策充分性；保留已确认英文正文及对应中文注释。
% \subsection{视觉动作模仿学习}
\subsection{Visuomotor Imitation Learning}
\label{sec:task}

% 我们研究有限、固定机器人示范下的离线视觉动作模仿学习，面向任务完成状态能够通过视觉历史评估的操作任务。
% 策略基于视觉、本体观测及可用的因果执行历史生成控制指令，部署期间不更新参数或查询专家。
% 训练以示范动作监督为基础，部署表现以有限时域内的任务完成概率衡量。
We study offline visuomotor imitation learning from a finite, fixed set of robot demonstrations,
for manipulation tasks whose completion can be assessed from visual history.
The policy generates control commands from visual and proprioceptive observations and available causal execution history,
without parameter updates or expert queries during deployment.
Training is based on demonstrated actions, whereas deployment performance is measured by the probability of task completion over a finite horizon.

% \subsubsection{系统与观测}
\subsubsection{System and Observations}
% 考虑由一个决策策略统一控制的机器人系统，通过各执行器的联合动作完成给定的操作任务。
% 系统与环境的联合状态 $s_t\in\mathcal S$ 在控制接口接收的指令 $u_t\in\mathcal A$ 的作用下演化，
% 并产生观测 $o_t$。在包含 $T$ 个控制步的有限时域内，该交互表示为如下部分可观测受控过程：
Consider a robotic system governed by a single decision policy that performs a given manipulation task
through the joint actions of its actuators.
The joint robot--environment state $s_t\in\mathcal S$ evolves under the command
$u_t\in\mathcal A$ and gives rise to observation $o_t$.
% 系统观测 $o_t=(I_t^{1:V},q_t)$ 由多视角图像 $I_t^{1:V}=(I_t^1,\ldots,I_t^V)$ 与本体信息 $q_t$ 组成，
% 其中 $V\ge1$ 为相机数量，$I_t^v$ 为第 $v$ 个视角的图像。
The observation $o_t=(I_t^{1:V},q_t)$ comprises multi-view images
$I_t^{1:V}=(I_t^1,\ldots,I_t^V)$ and proprioceptive information $q_t$,
where $V\ge1$ is the number of cameras and $I_t^v$ is the image from the $v$th view.
Over a finite horizon of $T$ control steps, this interaction is represented by the following partially observed controlled process:
\begin{align}
 s_0&\sim p_0,\nonumber\\
 s_{t+1}&\sim P(\cdot\mid s_t,u_t),\quad 0\le t<T,\nonumber\\
 o_t&\sim O(\cdot\mid s_t),\quad 0\le t\le T.
 \label{eq:interactionmodel}
\end{align}
% 其中，$\mathcal S$ 为状态空间，$\mathcal A\subseteq\R^{d_u}$ 为所选机器人控制接口规定的可行指令集合，
% $d_u$ 为指令维数，$p_0$ 为初始状态分布，$P$ 与 $O$ 分别为状态转移概率核和观测概率核。
Here, $\mathcal S$ is the state space, $\mathcal A\subseteq\R^{d_u}$ is the admissible command set
specified by the chosen robot control interface,
$d_u$ is the command dimension, $p_0$ is the initial-state distribution,
and $P$ and $O$ are the transition and observation kernels, respectively.

% \subsubsection{离线示范与策略}
\subsubsection{Offline Demonstrations and Policy}
% 策略从同一操作任务的固定示范集 $\mathcal D=\{\zeta_i\}_{i=1}^{N}$ 中学习，
% 第 $i$ 条示范轨迹 $\zeta_i$ 记录观测与示范控制指令：
The policy learns from a fixed demonstration set $\mathcal D=\{\zeta_i\}_{i=1}^{N}$
for the same manipulation task. The $i$th demonstration trajectory $\zeta_i$ records
observations and demonstrated control commands:
\begin{equation}
 \zeta_i=(o_0^i,u_0^{\mathrm E,i},o_1^i,\ldots,
 u_{T_i-1}^{\mathrm E,i},o_{T_i}^i).
 \label{eq:demonstrations}
\end{equation}
% 其中，$N$ 为有限的示范条数，$T_i\ge1$ 为第 $i$ 条轨迹的控制步数，上标 $\mathrm E$ 标识示范指令。
Here, $N$ is the finite number of demonstrations, $T_i\ge1$ is the number of control steps
in the $i$th trajectory, and the superscript $\mathrm E$ denotes demonstrated commands.

% 部署策略的信息由因果历史 $\mathcal H_t=(o_{0:t},u_{0:t-1},\mathcal R_{0:t})$ 表示，
% 其中 $\mathcal R_t$ 记录截至当前决策已获得的时间戳、计划及推理接收状态；无此记录时取空。
% 生成定义在动作集合 $\mathcal A$ 上的联合指令分布，并据此执行 $u_t\sim\pi_\theta(\cdot\mid\mathcal H_t)$。
% 其中，$\theta$ 为策略的学习参数，$t=0$ 时指令历史为空。
% 策略可采用单步动作预测、动作块预测或基础动作加残差的形式。
The information available to the deployed policy $\pi_\theta$ is represented by the causal history
$\mathcal H_t=(o_{0:t},u_{0:t-1},\mathcal R_{0:t})$,
where $\mathcal R_t$ records timestamps, plans, and inference-reception status available before the current decision; it is empty when no such records are used.
The policy produces a joint command distribution over the action set $\mathcal A$,
from which the executed command is drawn as $u_t\sim\pi_\theta(\cdot\mid\mathcal H_t)$.
Here, $\theta$ denotes the learned policy parameters;
the command history is empty at $t=0$.
The policy may use single-step action prediction, action-chunk prediction, or base actions augmented by learned residuals.

% \subsubsection{学习与任务目标}
\subsubsection{Learning and Task Objectives}
% 策略通过示范动作监督进行学习，动作监督目标按实际采样协议加权：
The policy learns from demonstrated actions, with the action-supervision objective expressed
as an empirical imitation risk weighted by the sampling protocol:
\begin{equation}
 \widehat R_{\mathrm{IL}}(\theta)
 =\sum_{i=1}^{N}\sum_{t=0}^{T_i-1}w_{it}
   \mathcal L_{\mathrm{IL}}\!\left(
   \pi_\theta;\mathcal H_t^i,\zeta_i\right).
 \label{eq:empiricalilrisk}
\end{equation}
% 其中，$\mathcal H_t^i$ 为示范 $\zeta_i$ 在时刻 $t$ 的记录历史，
% $\mathcal L_{\mathrm{IL}}$ 为衡量策略在该历史下的动作预测与对应示范动作差异的非负损失，
% $\zeta_i$ 提供相应的监督目标，包括动作块所需的未来动作序列。
Here, $\mathcal H_t^i$ is the recorded history at time $t$ in demonstration $\zeta_i$,
and $\mathcal L_{\mathrm{IL}}$ is a nonnegative loss measuring the discrepancy between
the policy's action predictions at that history and the corresponding demonstrated actions.
The trajectory $\zeta_i$ supplies the supervision targets, including future action sequences for action-chunk prediction.
% 其中 $w_{it}\ge0$ 为固定采样权重，总和为1；无效训练索引取零。
% 当前基础训练按可用帧索引采样，残差训练按构建后的缓存样本采样，不假定轨迹等权。
The fixed weights satisfy $w_{it}\ge0$ and $\sum_{i,t}w_{it}=1$, with zero weight on excluded indices.
The current base training samples available frame indices, whereas residual training samples constructed cache records; neither convention implies equal trajectory weighting.
% 具体动作预测形式、监督时间对齐、损失计算及辅助训练目标在方法部分给出。
The specific action-prediction forms, temporal alignment of supervision, loss computations,
and auxiliary training objectives are specified in the method section.

% 部署目标是在有限时域内完成任务。可测的任务完成函数 $\Phi:\mathcal S^{T+1}\to\{0,1\}$
% 对满足任务规定的终态或时序要求的状态轨迹取 $1$，否则取 $0$。
% 据此定义完成指示变量 $\Psi$ 与部署策略 $\pi$ 的任务完成概率 $J(\pi)$：
The deployment objective is to complete the task within the finite horizon.
The measurable task-completion function $\Phi:\mathcal S^{T+1}\to\{0,1\}$ equals $1$
for state trajectories satisfying the task's terminal or temporal requirements and $0$ otherwise.
This defines the completion indicator $\Psi$ and the task-completion probability $J(\pi)$ of the deployed policy $\pi$:
\begin{equation}
 \Psi=\Phi(s_{0:T}),\qquad
 J(\pi)=\Prob_\pi(\Psi=1).
 \label{eq:visualtask}
\end{equation}
% 其中，$\Prob_\pi$ 为系统模型与完整部署策略 $\pi$ 共同诱导的轨迹概率律。
% 示范监督用于学习策略，任务完成概率 $J(\pi)$ 用于评价其部署表现。
Here, $\Prob_\pi$ is the trajectory probability law induced jointly by the system model
and the complete deployed policy $\pi$.
Demonstration supervision trains the policy, whereas the task-completion probability
$J(\pi)$ evaluates its deployment performance.

% \subsubsection{视觉可评估性}
\subsubsection{Visual Assessability}
% 在固定相机配置与评价时域 $T$ 下，评价域
% $\Omega_{\mathrm{eval}}\subseteq\mathcal S^{T+1}\times\mathcal I^{T+1}$
% 由预先指定的、符合系统模型及可见性条件的物理轨迹与视觉历史对组成，
% 其中，$\mathcal I$ 为多视角图像元组 $I_t^{1:V}$ 的取值空间。
Under a fixed camera configuration and evaluation horizon $T$, the evaluation domain
$\Omega_{\mathrm{eval}}\subseteq\mathcal S^{T+1}\times\mathcal I^{T+1}$
consists of prespecified pairs of physical trajectories and visual histories
consistent with the system model and visibility conditions,
where $\mathcal I$ is the space of multi-view image tuples $I_t^{1:V}$.

% \begin{assumption}[视觉可评估性]
\begin{assumption}[Visual assessability]
\label{ass:visualassessability}
% 存在可测的视觉判别函数 $\mathcal V:\mathcal I^{T+1}\to\{0,1\}$，
% 使任意 $(s_{0:T},I_{0:T}^{1:V})\in\Omega_{\mathrm{eval}}$ 均满足
There exists a measurable visual decision function $\mathcal V:\mathcal I^{T+1}\to\{0,1\}$
such that, for every $(s_{0:T},I_{0:T}^{1:V})\in\Omega_{\mathrm{eval}}$,
\begin{equation}
 \Phi(s_{0:T})=\mathcal V(I_{0:T}^{1:V}).
 \label{eq:visassumption}
\end{equation}
\end{assumption}

% 假设~\ref{ass:visualassessability}仅限定评价域内完成事件的视觉可辨识性，
% 不要求完整状态可观测，也不意味着中间历史足以确定适当动作。
% 实际评价器误差的检验及域外轨迹的处理由实验协议另行规定。
Assumption~\ref{ass:visualassessability} only specifies visual identifiability of completion within the evaluation domain;
it neither requires full-state observability nor implies that intermediate histories suffice to determine appropriate actions.
The experimental protocol separately specifies assessment of practical evaluator errors and treatment of trajectories outside the domain.

% 核心观点：示范动作误差不是任务表现本身；任务改善取决于部署访问情形中的动作后果。
% 叙述逻辑：B1定义动作后果与性能差分；B2给出覆盖、敏感性条件下的动作误差界；
% B3用表示风险迁移的核心公式提出关系信息与执行更新两个问题；详细证明留附录，动作MSE诊断留D及附录。
% C、D继续研究表示与反馈的条件性作用，不预先声称语义降低C或反馈降低L。
% 边界：覆盖界只刻画动作风险迁移所需的条件；完整动作历史可能导致结构性支持失配。
% 不假定专家在离线数据之外可查询，也不由固定示范推出覆盖或随N变化的统计保证。
\subsection{From Imitation Error to Task Performance}
\label{sec:actionanalysis}
% 动作误差能否约束任务失败，取决于误差发生的情形及其后续任务后果；以下建立这一联系，并分析表示空间风险迁移的机会与限制。
Whether action error bounds task failure depends on where deviations occur and their subsequent task consequences.
We establish this connection and examine the opportunities and limitations of risk transfer through representations.

% \subsubsection{动作对任务的影响}
\subsubsection{Effects of Actions on Task Outcomes}
\label{sec:actionconsequences}
% 在同一执行情形下比较不同动作，随后均沿同一参考策略 $\pi_0$ 继续执行，
To compare actions in the same execution situation, we follow each action with the same reference policy $\pi_0$.
% 即可用最终失败概率的差异衡量当前动作对任务的影响。
The resulting difference in failure probability measures the effect of the current action on the task.
% 为保留任务时序与策略历史，引入分析用扩展状态 $x_t=(s_{0:t},\mathcal H_t)$。
To retain task timing and policy history, we introduce the augmented analysis state $x_t=(s_{0:t},\mathcal H_t)$.
% 后续分析限于固定的确定性策略，并简写 $\pi(x_t)=\pi(\mathcal H_t)$；
The following analysis considers fixed deterministic policies and uses the shorthand $\pi(x_t)=\pi(\mathcal H_t)$;
% 策略实际仅使用可获得的历史 $\mathcal H_t$。
the policy itself uses only the available history $\mathcal H_t$.

% 在控制步 $t$ 的执行情形 $x$ 下先执行动作 $a$、随后遵循 $\pi_0$ 时，
Let $Q_t^0(x,a)$ denote the conditional probability of task failure when action $a$ is executed
% 任务失败的条件概率记为 $Q_t^0(x,a)$。
in situation $x$ at control step $t$ and $\pi_0$ is followed thereafter.
% 执行参考动作时的失败概率为 $V_t^0(x)=Q_t^0(x,\pi_0(x))$，
The failure probability under the reference action is $V_t^0(x)=Q_t^0(x,\pi_0(x))$.
% 当前动作相对参考动作造成的失败概率差为
The change in failure probability relative to that action is
$A_t^0(x,a)=Q_t^0(x,a)-V_t^0(x)$.
% $A_t^0(x,a)>0$ 表示该动作增加失败概率，$A_t^0(x,a)<0$ 表示降低失败概率。
Thus, $A_t^0(x,a)>0$ indicates an increase in failure probability, whereas $A_t^0(x,a)<0$ indicates a decrease.

% \begin{assumption}[共同的比较过程]
\begin{assumption}[Common comparison process]
\label{ass:commonprocess}
% 比较策略共享初始分布、受控转移和任务完成函数。
The compared policies share the initial distribution, controlled transitions, and task-completion function.
% 包含运行记录时，共同转移也包括计算与记录过程，而不只是相同的机器人动力学。
When runtime records are included, the common transitions also cover the computation and recording process, not only the robot dynamics.
% 参考策略 $\pi_0$ 在比较所涉及的全部历史上有定义；
The reference policy $\pi_0$ is defined on every history involved in the comparison.
% 动作规则和上述条件期望可测，执行动作属于 $\mathcal A$。
The action rules and the conditional expectations above are measurable, and executed actions belong to $\mathcal A$.
\end{assumption}

% 以 $F(\pi)=1-J(\pi)$ 表示策略 $\pi$ 的任务失败概率。
Let $F(\pi)=1-J(\pi)$ denote the task-failure probability of policy $\pi$.
% 待评价策略诱导的时刻--历史分布 $\mu_\pi^t$ 描述其在控制步 $t$ 实际访问的 $x_t$，
The time--history distribution $\mu_\pi^t$ describes the augmented histories $x_t$ actually visited by the evaluated policy at control step $t$.
% 其等权时间混合记为 $\bar\mu_\pi=T^{-1}\sum_{t=0}^{T-1}\mu_\pi^t$，
Its uniform mixture over time is $\bar\mu_\pi=T^{-1}\sum_{t=0}^{T-1}\mu_\pi^t$,
% 其中混合的样本保留时间索引。
with each sample retaining its time index.
% 有限时域性能差分恒等式给出
The finite-horizon performance-difference identity gives
\begin{equation}
 F(\pi)-F(\pi_0)
 =T\,\E_{\bar\mu_\pi}\!\left[A_t^0(x,\pi(x))\right].
 \label{eq:pd}
\end{equation}
% 式\eqref{eq:pd}沿待评价策略 $\pi$ 实际访问的情形计算动作后果，
Equation~\eqref{eq:pd} evaluates action effects at the situations actually visited by the evaluated policy $\pi$.
% 这一标准性能差分关系\cite{cheng2018value,ross2014aggrevate}
This standard performance-difference relation~\cite{cheng2018value,ross2014aggrevate}
% 为后续分析提供共同尺度，证明见附录~\ref{app:actionproofs}。
provides a common measure for the subsequent analysis; a proof is given in Appendix~\ref{app:actionproofs}.

% \subsubsection{动作误差何时能够约束失败概率}
\subsubsection{When Action Error Bounds Failure Probability}
\label{sec:actionerrorbound}
% 以专家策略 $\pi_{\mathrm E}$ 为参考，下面考察平均动作误差在什么条件下能够约束任务失败概率。
Taking the expert policy $\pi_{\mathrm E}$ as the reference, we ask when mean action error can bound task-failure probability.
% 以 $\nu$ 表示评价动作误差的参考分布，其样本是控制时刻与完整执行历史组成的 $(t,x_t)$。
% 相应的平均动作平方误差为 $R_\nu(\pi)=\E_{(t,x)\sim\nu}\norm{\pi(x)-\pi_{\mathrm E}(x)}^2$。
Let $\nu$ be the reference distribution used to evaluate action error, with samples $(t,x_t)$ specifying a control step and its complete execution history.
% 在专家示范场景中，$\nu$ 可取为示范采集过程诱导的总体分布，而不是已采集有限样本的经验分布。
For expert demonstrations, $\nu$ may be the population distribution induced by the demonstration-collection process, rather than the empirical distribution of the collected samples.
The corresponding mean squared action error is $R_\nu(\pi)=\E_{(t,x)\sim\nu}\norm{\pi(x)-\pi_{\mathrm E}(x)}^2$.
% $\nu$ 描述动作误差在哪里被衡量，而 $\bar\mu_\pi$ 描述策略在部署时实际遇到的情况。
Whereas $\nu$ specifies where action error is evaluated, $\bar\mu_\pi$ describes the situations actually encountered during deployment.
% 以下两个条件足以建立任务表现界。
The following two conditions suffice to establish a task-performance bound.

% \emph{覆盖：}任意一类执行情况在部署中出现的概率，不超过参考分布中对应概率的 $C$ 倍：
\emph{Coverage:} For any class of execution situations, its probability during deployment must be at most $C$ times its probability under the reference distribution:
\[
 \bar\mu_\pi(\mathcal B)\le C\,\nu(\mathcal B)
 \quad\text{for every measurable event }\mathcal B.
\]
% 其中，$\mathcal B$ 为时刻--历史空间中的一类执行情况，$1\le C<\infty$ 为统一常数。
% 该条件既排除参考分布完全未覆盖的部署情况，也限制参考分布对其出现概率的低估。
Here, $\mathcal B$ is a set of time--history pairs representing a class of execution situations, and $1\le C<\infty$ is a uniform constant.
This condition excludes deployment situations absent from the reference distribution and limits how much their probabilities can be underestimated.

% \emph{敏感性：}在固定动作尺度下，偏离专家动作所引起的失败概率变化必须受动作距离约束：
\emph{Sensitivity:} For a fixed action scale, the change in failure probability caused by deviating from the expert action must be bounded by the action distance:
\[
 |A_t^{\mathrm E}(x,a)|
 \le L\norm{a-\pi_{\mathrm E}(x)}.
\]
% 其中，$A_t^{\mathrm E}$ 为以专家作为参考策略时的动作影响，$0\le L<\infty$ 为在所考察时刻、历史及动作上统一成立的常数。
Here, $A_t^{\mathrm E}$ is the action effect defined using the expert as the reference policy,
and $0\le L<\infty$ is a uniform constant over the times, histories, and actions considered.
% 在这两个条件下，性能差分恒等式给出如下失败概率上界（证明见附录~\ref{app:actionproofs}）：
Under these conditions, the performance-difference identity yields the following bound (Appendix~\ref{app:actionproofs}):
\begin{equation}
 F(\pi)\le F(\pi_{\mathrm E})+TL\sqrt{C R_\nu(\pi)}.
 \label{eq:successbound}
\end{equation}
% 这一推论沿用由部署模仿误差约束任务代价的思路\cite{ross2011}，在此通过覆盖条件和连续动作敏感性连接参考动作误差与失败概率。
This corollary follows the use of deployment imitation error to bound task cost~\cite{ross2011}, here relating reference action error to failure probability through coverage and continuous-action sensitivity.
% 其中C控制参考误差向部署误差的转换，L控制动作偏差向任务后果的转换。
The factor $C$ controls the transfer from reference to deployment action error, while $L$ controls the conversion of action deviations into task effects.
% 在其他量不变且条件成立时，减小R、C或L均收紧上界，但不直接给出不同策略实际成功率的排序。
With the other quantities fixed and the conditions satisfied, reducing $R_\nu$, $C$, or $L$ tightens the bound, without establishing an ordering of actual policy success rates.
% 有限示范上的训练误差与总体误差R的联系还需单独的泛化论证。
Connecting finite-demonstration training error to the population error $R_\nu$ requires a separate generalization argument.
% L可能很大或缺少有限统一界，识别关键情形本身并不改变物理系统的敏感性。
The constant $L$ may be large or lack a finite uniform bound; identifying critical situations does not itself change physical sensitivity.

% \subsubsection{表示层面的覆盖与条件误差差异}
\subsubsection{Representation-Level Coverage and Conditional Error Discrepancy}
% 完整历史覆盖可因小幅动作偏差失效。
Coverage over complete histories can fail after a small action deviation (Appendix~\ref{app:historycoverage}).
% 将历史编码为任务表示可以合并无关差异，但也可能合并需要不同动作的情形。
Encoding histories into a task representation can merge irrelevant differences, but may also merge situations requiring different actions.

% 固定策略、专家和同一个保留时间索引的可测因果编码器Z；mu_Z和nu_Z是部署与参考分布诱导的表示分布。
For a fixed policy, expert, and common measurable causal encoder $Z=\phi(t,\mathcal H_t)$ retaining the time index, let $\mu_Z$ and $\nu_Z$ be the representation distributions induced by $\bar\mu_\pi$ and $\nu$.
% 表示覆盖要求每类表示在部署中的出现概率不超过参考概率的C_phi倍。
Suppose $\mu_Z(\mathcal B)\le C_\phi\nu_Z(\mathcal B)$ for every measurable event $\mathcal B$, with $1\le C_\phi<\infty$.
% 同一表示下，两种分布仍可能包含不同历史；以D_phi衡量其条件平均动作平方误差的差异。
Even at the same representation, the two distributions may contain different histories.
Let $e(x)=\norm{\pi(x)-\pi_{\mathrm E}(x)}$ and define their conditional-error discrepancy as $D_\phi=\E_{z\sim\mu_Z}|m_\mu(z)-m_\nu(z)|$,
where $m_\mu(z)=\E_{\bar\mu_\pi}[e(x)^2\mid Z=z]$ and $m_\nu(z)=\E_\nu[e(x)^2\mid Z=z]$.
% 当动作平方误差在两分布下可积、条件均值存在可测版本且B2的敏感性条件成立时，有如下风险迁移及任务界。
With square-integrable action error under both distributions, measurable versions of these conditional means, and the sensitivity condition of Section~\ref{sec:actionerrorbound},
\begin{align}
 \E_{\bar\mu_\pi}e(x)^2
 &\le C_\phi R_\nu(\pi)+D_\phi,
 \label{eq:representationrisktransfer-main}\\
 F(\pi)&\le F(\pi_{\mathrm E})
       +TL\sqrt{C_\phi R_\nu(\pi)+D_\phi}.
 \label{eq:representationcoveragebound}
\end{align}
% 第一式按表示条件化后迁移动作风险，第二式接入B2的敏感性分析；详细证明见附录。
The first inequality transfers action risk by conditioning on the representation; the second applies the sensitivity argument of Section~\ref{sec:actionerrorbound}.
Appendix~\ref{app:representationcoverage} gives the detailed proof.
% 该界同时取决于覆盖加权的参考风险与条件误差差异，覆盖常数较小本身不足以收紧上界。
The bound depends jointly on the coverage-weighted reference risk and the conditional error discrepancy; a smaller coverage factor alone need not tighten it.
% 保留RGB后附加确定性语义不自动改善覆盖。
Appending a deterministic semantic map to retained RGB does not by itself improve coverage.
% 对固定策略、专家与编码器，D_phi比较两分布的条件平均平方动作误差，而非任务信息损失；合并不同所需动作的历史也可能使其为零。
For the fixed policy, expert, and encoder, $D_\phi$ compares conditional mean squared action errors across distributions, rather than measuring lost task information.
It can vanish even when histories requiring different actions are merged, so a zero discrepancy does not establish information sufficiency.
% C直接以任务后果评价信息保留，与风险迁移分析互补；不界定D_phi，也不依赖上述覆盖条件。
Section~\ref{sec:representationanalysis} complements risk-transfer analysis by assessing information preservation directly through task consequences; it neither bounds $D_\phi$ nor requires the preceding coverage conditions.

% 核心观点：表示应保留区分动作任务后果所需的关系与因果执行上下文，
% 而不只是使专家动作MSE的上界更小。
% 叙述逻辑：信息保留/决策利用分解 -> 关系保留的条件性界 -> 语义与关系监督的对应。
% 正文集中比较条件与适用域；标量实例和伪标签误差分解留附录，检验协议见实验。
% 边界：语义是引导关系表示的手段，不另设动作预测风险或泛化收益命题；
% 关系界约束可达决策能力，不证明特定架构优于RGB基线。
% 任务后果代理仅为分析见证函数；几何辅助损失不等于价值监督或策略改进证明。
\subsection{Semantic Inputs and Task-Supervised Representations}
\label{sec:representationanalysis}
% 任务相关表示需要保留决定动作后果的关系与上下文；以下先区分信息损失与决策实现误差，再分析关系估计误差如何约束表示损失。
Task-relevant representations must preserve the relations and context needed to distinguish action outcomes.
We first separate information loss from decision-realization error, then bound representation loss through relation-estimation error.

% \subsubsection{表示信息与决策能力}
\subsubsection{Represented Information and Decision Capability}
% 固定参考策略 $\pi_0$ 及一个时刻--历史评价分布 $\mu$。
Fix a reference policy $\pi_0$ and a time--history evaluation distribution $\mu$.
% 动作规则 $a$ 对应的平均失败概率定义为
Define the mean failure probability associated with an action rule $a$ as
$\mathcal L_\mu(a)=\E_\mu[Q_t^0(x,a)]$.
% 它衡量替换当前动作后沿 $\pi_0$ 继续执行的平均失败概率。
This quantity measures the mean failure probability after replacing the current action and following $\pi_0$ thereafter.

% 含时间索引的完整因果历史生成信息 $\mathcal F_h=\sigma(\mathcal H_t)$。
The complete causal history, including the time index, generates the information $\mathcal F_h=\sigma(\mathcal H_t)$.
% 编码表示 $Z=\phi(\mathcal H_t)$ 生成的信息为
% Z在此表示目标指令的完整决策输入：除学习特征外，还包含目标时刻/槽位、所需计划缓存及因果调度记录。
The representation $Z=\phi(\mathcal H_t)$ denotes the complete input determining the target command: learned features together with the target time and slot, required cached plan information, and causal scheduling records.
% 这些量均由含运行记录的因果历史确定；确定性读出包括槽位选择及所用后处理。
These quantities are determined by the causal history including runtime records; the deterministic readout includes slot selection and any postprocessing used.
% Z止于待评价决策读出之前；反馈可条件于已有基础计划，但不将其自身的最终修正动作作为输入。
The representation is taken before the evaluated decision readout: feedback may condition on an existing base plan, but not on its own final corrected command.
It generates
$\mathcal U=\sigma(Z)\subseteq\mathcal F_h$.
% 基础策略与反馈分支的不同编码分别作为 $Z$ 评价。
The encodings of the base policy and the feedback branch are evaluated separately as $Z$.
% 对任意信息集合 $\mathcal F$，以 $\mathfrak A(\mathcal F)$ 表示所有
For any information set $\mathcal F$, let $\mathfrak A(\mathcal F)$ denote all
% $\mathcal F$-可测且取值于 $\mathcal A$ 的动作规则，
$\mathcal F$-measurable action rules taking values in $\mathcal A$.
% 仅使用这些信息所能达到的平均失败概率下确界为
The infimum of the mean failure probability attainable using only this information is
\begin{equation}
 \mathcal L_\mu^*(\mathcal F)
 =\inf_{a\in\mathfrak A(\mathcal F)}\mathcal L_\mu(a).
 \label{eq:taskinformationrisk}
\end{equation}
% 其中sigma表示记录可辨别的信息，F可测规则仅依赖该信息。
Here, $\sigma(\cdot)$ denotes the information generated by the records, and an $\mathcal F$-measurable rule uses only that information.
% 各类动作规则在相同动作集合下比较，当前动作之后均沿同一参考策略执行。
The action-rule classes are compared over the same action set, with the same reference policy followed after the current action.

% 基于表示的实际动作规则 $\pi_Z$ 同时受到编码信息与决策机制的限制。
An implemented action rule $\pi_Z\in\mathfrak A(\mathcal U)$ is constrained both by the encoded information and by its decision mechanism.
% 将编码丢失信息所造成的最优平均失败概率差记为表示损失
Define the gap in optimal mean failure probability due to information lost in encoding as the representation loss
$\Delta_{\mathrm{rep}}(\mathcal U)
=\mathcal L_\mu^*(\mathcal U)-\mathcal L_\mu^*(\mathcal F_h)$,
% 将实际决策相对同一表示下最优规则的差距记为决策实现误差
and the gap between the implemented rule and the optimum under the same representation as the decision-realization error
$\epsilon_{\mathrm{dec}}(\pi_Z)
% =\mathcal L_\mu(\pi_Z)-\mathcal L_\mu^*(\mathcal U)$，有
=\mathcal L_\mu(\pi_Z)-\mathcal L_\mu^*(\mathcal U)$. Then
\begin{equation}
 \mathcal L_\mu(\pi_Z)-\mathcal L_\mu^*(\mathcal F_h)
 =\Delta_{\mathrm{rep}}(\mathcal U)+\epsilon_{\mathrm{dec}}(\pi_Z).
 \label{eq:taskrepresentationdecomposition}
\end{equation}
% 该分解区分\emph{信息是否被保留}与\emph{保留的信息是否被有效利用}。
This decomposition distinguishes \emph{whether information is preserved} from \emph{whether the preserved information is effectively used}.
% 两项均以沿参考策略继续执行的平均失败概率计量且非负；$\epsilon_{\mathrm{dec}}$ 包含模型容量、残差限幅、
Both terms are nonnegative and measured in mean failure probability under reference-policy continuation. The term $\epsilon_{\mathrm{dec}}$ includes limitations due to model capacity, residual bounds,
% 有限数据与优化，以及动作监督和任务目标之间的差异。
finite data and optimization, and the mismatch between action supervision and the task objective.

% \subsubsection{关系误差何时影响任务决策}
\subsubsection{When Relation Errors Affect Task Decisions}
% 关系估计的准确性需要通过动作后果与任务决策联系起来。
The accuracy of relation estimates must be connected to task decisions through action outcomes.
% 给定可获得的历史，对未知物理状态取条件平均后的失败概率记为
Given the available history, denote the failure probability conditionally averaged over unknown physical states by
$\overline Q_\mu(\mathcal H_t,a)=\E_\mu[Q_t^0(x,a)\mid\mathcal F_h]$,
% 其中 $a$ 为固定的待比较动作，后续均沿参考策略执行。
where $a$ is the fixed action being compared, followed in each case by the reference policy.

% 关系变量 $G\in\R^{d_G}$ 描述相关对象的几何或语义关系，
The relation variable $G\in\R^{d_G}$ describes geometric or semantic relations among relevant objects,
% 因果上下文 $\chi$ 保留决策所需的本体信息与执行证据。
and the causal context $\chi$ retains proprioceptive information and execution evidence needed for decisions.
% chi的组成须在任务实例中预先限定，不能通过把完整历史放入chi使假设恒成立；未解释因素由epsilon_rel承担。
The contents of $\chi$ must be specified for the task instance, rather than set to the entire history to make the assumption automatic; unexplained effects remain in $\epsilon_{\mathrm{rel}}$.
% 表示 $Z$ 可读出关系估计 $\widehat G\in\R^{d_G}$ 与上下文 $\chi$，
The representation $Z$ admits readouts of a relation estimate $\widehat G\in\R^{d_G}$ and context $\chi$,
% 关系误差为 $\delta_G=\E_\mu\norm{G-\widehat G}$。
with relation error $\delta_G=\E_\mu\norm{G-\widehat G}$.
% 以下用分析函数 $f(G,\chi,a)$ 近似上述条件平均失败概率，
We use an analytical function $f(G,\chi,a)$ to approximate the conditional mean failure probability above.
% 其中d_G为关系维数。
Here, $d_G$ is the relation dimension.

% \begin{assumption}[关系条件下的任务后果近似]
\begin{assumption}[Relation-based approximation of action outcomes]
\label{ass:taskrelations}
% 在固定分布 $\mu$ 下，存在可测函数 $f$ 及常数 $\epsilon_{\mathrm{rel}},K_G\ge0$，使
Under a fixed distribution $\mu$, there exist a measurable function $f$ and constants $\epsilon_{\mathrm{rel}},K_G\ge0$ such that
\begin{align}
 \E_\mu\sup_{a\in\mathcal A}
 |\overline Q_\mu(\mathcal H_t,a)-f(G,\chi,a)|
 &\le\epsilon_{\mathrm{rel}},\nonumber\\
 |f(G,\chi,a)-f(\widehat G,\chi,a)|
 &\le K_G\norm{G-\widehat G},\quad a\in\mathcal A .
 \label{eq:taskrelationsassumption}
\end{align}
% 第二个条件在真实与估计关系的联合支持上几乎处处成立，
The second condition holds almost everywhere on the joint support of the true and estimated relations,
% 且 $\delta_G<\infty$；条件期望、上确界和动作选择满足附录~\ref{app:representationregularity}的正则条件。
and $\delta_G<\infty$. The conditional expectations, suprema, and action selections satisfy the regularity conditions in Appendix~\ref{app:representationregularity}.
\end{assumption}
% 其中，$\epsilon_{\mathrm{rel}}$ 为关系和上下文不足以解释任务后果的近似误差，
Here, $\epsilon_{\mathrm{rel}}$ is the approximation error due to task effects not captured by the relations and context,
% $K_G$ 为该后果对关系误差的敏感性常数，$\delta_G$ 为关系估计误差。
$K_G$ is the sensitivity constant with respect to relation errors, and $\delta_G$ is the relation-estimation error.
% 例如，相同几何下的接触或滑脱差异须由上下文保留，否则计入近似误差。
For example, contact and slip at the same geometry must be distinguished by the context or accounted for in the approximation error.

% \begin{proposition}[任务后果意义下的表示损失]
\begin{proposition}[Representation loss in task outcomes]
\label{prop:taskrepresentation}
% 在假设~\ref{ass:taskrelations}下，
Under Assumption~\ref{ass:taskrelations},
\begin{align}
 \Delta_{\mathrm{rep}}(\mathcal U)
 &\le 2(\epsilon_{\mathrm{rel}}+K_G\delta_G),\nonumber\\
 \mathcal L_\mu(\pi_Z)-\mathcal L_\mu^*(\mathcal F_h)
 &\le 2(\epsilon_{\mathrm{rel}}+K_G\delta_G)
       +\epsilon_{\mathrm{dec}}(\pi_Z).
 \label{eq:taskrelationbound}
\end{align}
\end{proposition}
% 证明见附录~\ref{app:representationproofs}。
The proof is given in Appendix~\ref{app:representationproofs}.
% RGB和增强编码器比较必须固定mu、pi0、动作集合与尺度、G定义与尺度、上下文及执行输入。
Comparisons between RGB and enhanced encoders use the same $\mu$, reference continuation $\pi_0$, action set and scale, relation definition and normalization, and context and execution inputs.
% 分别训练的表示并不必然信息嵌套，须分别检验假设与误差；仅当共享近似误差和敏感性上界时才能把更小delta解释为更紧上界。
Separately trained representations need not have nested information sets: the assumption and errors are evaluated for each, and smaller $\delta_G$ gives a tighter common bound only with shared bounds on $\epsilon_{\mathrm{rel}}$ and $K_G$.
% 该比较不从关系误差大小直接推出实际策略排序。
Relation-readout accuracy alone does not rank their implemented policies.
% 附录的理想化末步纠偏例子展示非零、有限关系敏感性下的信息价值，并非实际接触动力学模型。
An idealized final-step correction in Appendix~\ref{app:relationexample} illustrates the value of relations under nonzero, finite sensitivity; it is not a model of the actual contact dynamics.
% 关系因遮挡而未定义时，关系界仅适用于有效条件分布，无效情形仍计入总体任务结果。
When occlusion leaves relations undefined, the relation bound applies only to the valid conditional distribution, while invalid cases remain part of the overall task results.

% 关系读出与伪标签误差均需独立核查，不能以训练拟合代替理论中的评价误差。
Relation-readout and pseudo-label errors require independent assessment (Appendix~\ref{app:pseudolabel}); training fit cannot substitute for the evaluation errors used in the analysis.

% 上述结果表明，在相应条件下，准确保留任务相关关系与上下文可以约束表示损失；实际决策表现还取决于策略对这些信息的利用。
Under the stated conditions, accurately preserving task-relevant relations and context in $Z$ bounds representation loss; actual decision performance also depends on how the policy uses this information.
% 下一节进一步讨论执行中新增证据的利用。
The next section examines the use of additional evidence acquired during execution.

% 核心观点：反馈只有把新增的任务相关证据转化为及时、有效的动作改变，才改善任务完成概率。
% 叙述逻辑：执行变化提出新增证据的需求 -> D1刻画信息的潜在价值 -> D2调用C的信息保留/利用分解 ->
% D3分析实际采用，并调用B1性能差分连接闭环任务结果；不是为历史与残差另建独立分析。
% C到D：关系界须对反馈联合表示及执行时刻的关系成立，不沿用基础计划时刻的几何误差。
% 边界：不把GRU/MLP误差差值当成理想历史价值，不把残差幅度或动作MSE改善当成任务改进。
% \subsection{历史、修正与任务收益}
\subsection{History, Corrections, and Task Benefit}
\label{sec:feedbackanalysis}
% 新增执行证据的收益取决于原先可用的信息，以及所生成的修正是否被实际采用。
The benefit of additional execution evidence depends on the information already available and on whether the resulting correction is actually adopted.

% \subsubsection{新观测与历史的决策价值}
\subsubsection{Decision Value of New Observations and History}
% 以下以动作块执行为具体场景，比较已有计划给出的基础动作与利用新增执行信息生成的修正动作。
We take action-chunk execution as the concrete setting, comparing a base action supplied by an existing plan with a corrected action generated using additional execution information.
% 在固定评价分布 $\mu$ 上，以 $t$ 表示目标执行时刻，将该控制步的动作位置称为执行槽位。
Under a fixed evaluation distribution $\mu$, let $t$ denote the target execution time, and call the action position at that control step an execution slot.
% 基础策略为该槽位选中的动作记为
Denote the action selected by the base policy for this slot by
$B=\pi_0(x)$.
% 这里 $\pi_0$ 为包含计划更新、槽位选择及预先规定回退流程的完整基础执行规则。
Here, $\pi_0$ is the complete base execution rule, including plan updates, slot selection, and predefined fallback procedures.
% 本节参考策略保留相同的基础与反馈推理、接收和调度流程，仅在执行输出端屏蔽残差。
The reference policy retains the same base and feedback inference, reception, and scheduling procedures, but masks the residual at the execution output.
% 计划信息 $\mathcal F_p$ 包含生成和选择该动作所使用的观察与计划上下文，
The planning information $\mathcal F_p$ contains the observations and plan context used to generate and select this action,
% 以及触发回退所需的运行记录，使 $B$ 对 $\mathcal F_p$ 可测。
together with runtime records needed to trigger fallback, so that $B$ is $\mathcal F_p$-measurable.
% 当前信息 $\mathcal F_c$ 在保留计划信息的基础上加入时刻 $t$ 最新可用的观测和本体信息；
The current information $\mathcal F_c$ retains the planning information and adds the latest observations and proprioception available at time $t$.
% 完整因果历史 $\mathcal F_h$ 进一步包含此前的执行记录，
The complete causal history $\mathcal F_h$ additionally contains previous execution records,
% 满足 $\mathcal F_p\subseteq\mathcal F_c\subseteq\mathcal F_h$。
with $\mathcal F_p\subseteq\mathcal F_c\subseteq\mathcal F_h$.
% 运行时信息仅包含已获得的记录及时间戳。
Runtime information includes only records and timestamps already available.

% 新观测和额外历史分别使最优平均失败概率降低的幅度为
The reductions in optimal mean failure probability due to new observations and additional history are, respectively,
\begin{align}
 \Delta_{\mathrm{new}}
 &=\mathcal L_\mu^*(\mathcal F_p)-\mathcal L_\mu^*(\mathcal F_c),\nonumber\\
 \Delta_{\mathrm{hist}}
 &=\mathcal L_\mu^*(\mathcal F_c)-\mathcal L_\mu^*(\mathcal F_h).
 \label{eq:taskinformationgain}
\end{align}
% 信息嵌套使两项均非负，因为较丰富的信息允许复用原动作规则。
Both terms are nonnegative by information inclusion: a rule using less information remains available when more information is provided.
% 一个严格正历史价值的充分情形是：相同当前观测和计划下，两种等概率模式要求相反动作，历史识别模式，正确与错误动作的失败概率差为c>0。
% 可选动作仅为两个相反动作；相同当前信息指全部F_c，包括本体、计划和运行记录；评价分布仅在末次决策上。
For a final-decision distribution $\mu$ in a two-action setting, consider two equally likely modes with identical $\mathcal F_c$, including observations, proprioception, plan context, and runtime records.
If history distinguishes the modes and they require opposite actions, with an incorrect choice increasing failure probability by $0<c\le1$, then $\Delta_{\mathrm{hist}}=c/2$ (Appendix~\ref{app:historyvalue}).
% 该例不是声称所有历史都有价值；若当前信息已决定最优动作，历史价值可为零。
If current information already determines an optimal action, the history value can instead be zero.
% 有限动作--响应历史的编码可能丢失决策所需信息，该损失由后续表示损失项刻画。
Encoding a finite action--response history may discard decision-relevant information; the resulting loss is accounted for by the representation term below.

% \subsubsection{从信息价值到实际修正}
\subsubsection{From Information Value to Implemented Corrections}
% 新信息的潜在价值不等于候选修正的实际收益；沿用C的分解，还须扣除反馈表示的信息损失与决策实现误差。
The potential value of new information is not yet the benefit of a candidate correction.
Applying the decomposition in~\eqref{eq:taskrepresentationdecomposition} to the feedback representation also accounts for information loss and decision-realization error.
% 实际反馈表示 $Z_t^{\mathrm{fb}}$ 编码基础计划、缓存查询、RGB、因果历史及候选选择所需运行记录，
The implemented feedback representation $Z_t^{\mathrm{fb}}$ encodes the base plan, cached queries, RGB inputs, causal history, and runtime records required for candidate selection,
% 生成信息 $\mathcal U=\sigma(Z_t^{\mathrm{fb}})\subseteq\mathcal F_h$。
generating $\mathcal U=\sigma(Z_t^{\mathrm{fb}})\subseteq\mathcal F_h$.
% 在基础计划有效时，基础查询对应计划观测时刻 $\tau_p$，并在计划内保持不变；
When the base plan is valid, its queries correspond to the planning observation time $\tau_p$ and remain unchanged within the plan.
% 候选修正使用最新观测时刻 $\tau_o$ 的图像，作用于槽位 $t$，其中 $\tau_p\le\tau_o\le t$。
Candidate corrections use images from the latest observation time $\tau_o$ and target slot $t$, where $\tau_p\le\tau_o\le t$.
% 编码过程中的信息丢失与观测延迟共同限制该表示所保留的信息。
Information loss during encoding and observation delay jointly limit the information retained by this representation.

% 由该表示生成候选指令 $\widehat u=B+r\in\mathcal A$，
The representation produces a candidate command $\widehat u=B+r\in\mathcal A$,
% 其中 $r$ 为包含限幅等后处理的目标槽位残差。
where $r$ is the residual for the target slot after postprocessing, including magnitude bounding.
% 基础动作的决策实现差距为
The decision-realization gap of the base action is
$\epsilon_{\mathrm{base}}=\mathcal L_\mu(B)-\mathcal L_\mu^*(\mathcal F_p)$.

% \begin{proposition}[候选反馈的任务收益分解]
\begin{proposition}[Decomposition of candidate feedback benefit]
\label{prop:taskfeedback}
% 在同一 $\mu$、$\pi_0$ 和动作集合下，将采用候选修正后平均失败概率的下降称为候选反馈收益，满足
Under the same $\mu$, $\pi_0$, and action set, define the candidate feedback benefit as the reduction in mean failure probability from applying the candidate correction. It satisfies
\begin{align}
 \Gamma_\mu
 &:=\mathcal L_\mu(B)-\mathcal L_\mu(\widehat u)\nonumber\\
 &=\epsilon_{\mathrm{base}}+\Delta_{\mathrm{new}}+\Delta_{\mathrm{hist}}
   -\Delta_{\mathrm{rep}}(\mathcal U)-\epsilon_{\mathrm{dec}}(\widehat u).
 \label{eq:taskfeedbackdecomposition}
\end{align}
\end{proposition}
% 证明见附录~\ref{app:feedbackproofs}。
The proof is given in Appendix~\ref{app:feedbackproofs}.
% 决策实现差距包含残差幅度/容量限制、学习误差及监督目标差异。
The decision-realization gap includes residual magnitude and capacity limits, learning error, and the mismatch between action supervision and task consequences.
% 同分布同槽位的残差动作误差恒等式仅诊断修正方向与幅度，不能替代这里的任务收益。
The matched-distribution, matched-slot action-error identity~\eqref{eq:residualgain} diagnoses correction direction and magnitude, but does not replace this task-benefit criterion.

% \subsubsection{时序对齐与闭环任务收益}
\subsubsection{Temporal Alignment and Closed-Loop Task Benefit}
% 候选的决策优势只有通过实际采用才影响轨迹；以下结合采用规则与B1性能差分。
A candidate's decision advantage affects the trajectory only through its actual adoption.
We connect the adoption rule to task performance using the performance-difference identity~\eqref{eq:pd}.
% r是执行机制实际选中并限幅的目标槽位候选，无候选取0；候选必须及时到达、未过期且匹配有效计划与槽位。
Let $r$ be the bounded target-slot candidate actually selected by the deterministic execution rule, or zero if none is available.
A correction must arrive in time, remain unexpired, and match the valid plan and execution slot.
% 修正采用指示量 $\Lambda\in\{0,1\}$ 在满足这些条件时取 $1$，否则取 $0$。
The correction-adoption indicator $\Lambda\in\{0,1\}$ equals $1$ when these conditions hold and $0$ otherwise.
% 无有效修正时执行 $B$；基础计划也不可用时，$B$ 为同一历史下预先规定的回退指令，
Without a valid correction, $B$ is executed. If the base plan is also unavailable, $B$ is the predefined fallback command for the same history,
% 此时置 $\Lambda=0$，并将候选残差约定为 $r=0$。
with $\Lambda=0$ and the candidate residual set to $r=0$.
% 实际执行指令为 $u_{\mathrm{exec}}=B+\Lambda r$，
The executed command is $u_{\mathrm{exec}}=B+\Lambda r$.
% 在执行情形 $x$ 下，候选修正相对基础动作降低失败概率的幅度为
In situation $x$, the candidate's reduction in failure probability relative to the base action is
$g_t(x)=Q_t^0(x,B)-Q_t^0(x,\widehat u)$.
% 由于任务后果是失败概率，$g_t(x)\in[-1,1]$。
Because the compared outcomes are failure probabilities, $g_t(x)\in[-1,1]$.
% 有效收益为采用后的平均任务收益，取决于哪些修正及时采用，而非仅总体采用率。
The effective benefit is $\Gamma_{\mathrm{exec},\mu}=\E_\mu[\Lambda g_t(x)]$;
it depends on which corrections are adopted in time, not merely the overall adoption rate.

% 参照保留相同计算与接收流程，以符合共同过程条件；直接关闭反馈计算可能改变基础计划时序，是不同的系统级比较。
The output-masked reference retains the same computation and reception process, as required by Assumption~\ref{ass:commonprocess}.
Skipping feedback computation can change base-plan timing under the shared-worker schedule and therefore defines a different system-level comparison.

% \begin{proposition}[有效反馈与任务完成]
\begin{proposition}[Effective feedback and task completion]
\label{prop:executedtaskgain}
% 在假设~\ref{ass:commonprocess}下，
Under Assumption~\ref{ass:commonprocess},
% 若完整部署策略 $\pi$ 在其访问历史上执行上述 $u_{\mathrm{exec}}$，
if the complete deployed policy $\pi$ executes the above $u_{\mathrm{exec}}$ on its visited histories,
% 而参考策略在相同历史上执行 $B=\pi_0(x)$，则
while the reference policy executes $B=\pi_0(x)$ on the same histories, then
\begin{equation}
 J(\pi)-J(\pi_0)
 =T\,\E_{\bar\mu_\pi}[\Lambda g_t(x)].
 \label{eq:closedlooptaskgain}
\end{equation}
% 因此，实际访问分布上的有效任务收益为正，当且仅当任务完成概率高于参考策略。
Thus, the effective task benefit under the actual visitation distribution is positive if and only if task-completion probability exceeds that of the reference policy.
\end{proposition}

% 进一步，若 $Z_t^{\mathrm{fb}}$ 可读出执行时刻关系 $G_t$ 的估计 $\widehat G_t$ 与上下文 $\chi_t$，
Furthermore, if $Z_t^{\mathrm{fb}}$ admits readouts of an estimate $\widehat G_t$ of the execution-time relations $G_t$ and context $\chi_t$,
% 且这些量在 $\bar\mu_\pi$ 上满足假设~\ref{ass:taskrelations}，则可应用C的关系界。
and these quantities satisfy Assumption~\ref{ass:taskrelations} under $\bar\mu_\pi$, the relation bound in Section~\ref{sec:representationanalysis} applies.
% 此处 $\delta_G=\E_{\bar\mu_\pi}\norm{G_t-\widehat G_t}$ 及其余关系常数均针对反馈表示与执行时刻；
Here, $\delta_G=\E_{\bar\mu_\pi}\norm{G_t-\widehat G_t}$ and all other relation constants concern the feedback representation and execution time.
% 计划时刻 $\tau_p$ 的查询精度不能替代它们（附录~\ref{app:feedbackrelationtime}）。
Query accuracy at the planning time $\tau_p$ cannot replace these quantities (Appendix~\ref{app:feedbackrelationtime}).
% 以 $p_{\mathrm{miss}}=\Prob_{\bar\mu_\pi}(\Lambda=0)$ 表示未采用概率，结合时效界得到
Let $p_{\mathrm{miss}}=\Prob_{\bar\mu_\pi}(\Lambda=0)$ denote the probability of nonadoption. Combining this with the timing bound gives
\begin{equation}
 \begin{aligned}
 J(\pi)-J(\pi_0)\ge T\big[&
 \epsilon_{\mathrm{base}}+\Delta_{\mathrm{new}}+\Delta_{\mathrm{hist}}\\
 &-2(\epsilon_{\mathrm{rel}}+K_G\delta_G)\\
 &-\epsilon_{\mathrm{dec}}(\widehat u)-p_{\mathrm{miss}}\big].
 \end{aligned}
 \label{eq:endtoendtaskbound}
\end{equation}
% 右端为正时，信息价值及基础决策的改进空间超过表示、实现和时效损失，
A positive right-hand side means that information value and the room for improving the base decision exceed representation, realization, and timing losses,
% 从而给出任务完成概率提高的充分条件。
providing a sufficient condition for increased task-completion probability.
% 所有量在同一部署分布评价；这是固定基础策略上的反馈增益，不能将不同策略分布上的局部收益相加。
All terms are evaluated under the same deployment distribution for a fixed base policy.
If $\pi_0$ already includes semantics, this is a feedback-increment guarantee, not the full method's gain over RGB; local gains under different policy distributions cannot simply be added.
% 若关系条件仅在可见等有效子域成立，关系界只适用于该条件分布；
If the relation conditions hold only on a valid subdomain, such as visible cases, the relation bound applies only to that conditional distribution.
% 总体收益还须计入无效子域的贡献，不能直接由式\eqref{eq:endtoendtaskbound}获得保证。
The overall benefit must also account for the invalid subdomain and cannot be guaranteed directly by~\eqref{eq:endtoendtaskbound}.

% 动作与几何监督本身不确保该充分条件成立；回放迁移与动作诊断见附录。
Action and geometric supervision alone do not establish this sufficient condition; replay-transfer bounds and action diagnostics are provided in Appendix~\ref{app:feedbackproofs}.

\subsection{Testable Implications}
分析给出三个需要分别检验的环节，而非用单一成功率验证全部假设。
首先，在统一动作尺度、执行槽位和独立轨迹上比较动作风险，
并区分接触等关键事件附近与普通状态的误差，以检验平均误差的解释边界。
其次，语义与查询的收益应同时体现为相关几何可读出、策略确实利用该表示以及闭环行为改善；
仅降低辅助损失不足以完成这一论证。
最后，历史反馈应在同一冻结基础策略上检验残差与误差的对齐、历史的增量贡献和延迟鲁棒性。
上述实验可支持机制解释，但不能从有限样本证明全状态空间的覆盖或稳定性假设。
\draftnote{待补：各命题对应的实验编号、预注册比较项及判定标准；未知的覆盖常数和敏感性常数不填入推测数值。}
